\chapter{Trabalhos Relacionados}
\label{cap:relacionados}

Este capítulo apresenta trabalhos que aplicaram verificação formal a sistemas
embarcados e a software de veículos autônomos, posicionando esta pesquisa em
relação ao estado da arte.

\section{Protocolo de Levantamento}
\label{sec:rel-protocolo}

O levantamento foi exploratório, e não uma revisão sistemática: o objetivo foi
posicionar a pesquisa, não mapear exaustivamente o campo. As bases consultadas
foram IEEE Xplore, ACM Digital Library e Google Scholar --- esta usada sobretudo
para rastrear citações a partir dos trabalhos já identificados. As buscas
combinaram termos de três eixos: a técnica (\textit{bounded model checking},
\textit{formal verification}), o alvo (\textit{embedded firmware}, \textit{UAV},
\textit{autopilot}) e a ferramenta (ESBMC, CBMC, Frama-C), privilegiando
publicações a partir de 2019, exceto os trabalhos seminais da área.

Foram incluídos trabalhos que aplicam verificação formal a código C/C++, que
tratam de sistemas embarcados críticos ou de veículos autônomos, ou que levantam
o estado da arte nesses domínios. Excluíram-se os restritos a verificação em
nível de modelo sem correspondência com código executável e os que tratam de
segurança de UAVs sob perspectiva exclusivamente de rede ou ataque físico, sem
relação com o firmware embarcado.
%% LC (acknowledged): protocolo exploratorio; nao ha registro de strings/datas/
%% contagens -- retirada a promessa de reproduzir a busca.

\section{Verificação Formal de Sistemas Embarcados}
\label{sec:rel-embarcados}

Ferramentas de BMC como CBMC~\cite{Kroening2014} e ESBMC~\cite{Menezes2024}
permitem analisar programas C/C++ pela exploração simbólica de seus caminhos de
execução. O CBMC foi uma das primeiras a aplicar BMC diretamente a programas C,
adotando uma abordagem baseada em SAT; o ESBMC adota uma abordagem baseada em
SMT~\cite{deMoura2008}, que incorpora teorias como vetores de bits,
\textit{arrays} e ponto flutuante. Aplicações recentes demonstram o uso dessas
ferramentas na identificação de defeitos em projetos de código aberto de grande
escala~\cite{Sousa2026}. As duas são usadas nesta pesquisa, a segunda na
validação cruzada.

\textit{Surveys} recentes~\cite{Beyer2022, Beyer2024} indicam que boa parte dos
trabalhos de verificação ainda se concentra em programas sintéticos ou
\textit{benchmarks} padronizados, e não em sistemas reais de produção --- lacuna
que motiva diretamente este trabalho. Cabe a ressalva de que essa observação
decorre de um levantamento exploratório, não de uma medição sistemática.
%% LC (resolvido): afirmacao de lacuna atenuada -- nao atribuida aos relatorios
%% de competicao, que nao a sustentam.

\section{Aplicações Práticas de Verificação Formal}
\label{sec:rel-praticas}

Sousa et al.~\cite{Sousa2026} aplicaram o ESBMC a projetos \textit{open-source}
em C de grande escala, com uma metodologia sistemática de pré-processamento
automatizado e priorização por função para guiar o verificador nas partes mais
críticas, identificando defeitos previamente desconhecidos. É a abordagem mais
próxima desta pesquisa; a diferença é que os projetos de Sousa et al. são
aplicações genéricas, enquanto este estudo foca em firmware de controle de voo
com restrições de tempo real --- tratadas aqui pela seleção de módulos
verificáveis isoladamente, com \textit{stubs} que abstraem dependências de
hardware e temporização, sem comprometer a análise das propriedades de
\textit{safety} internas a cada módulo.
%% LC (resolvido): Sousa2023 era o preprint de Sousa2026 -- citacoes unificadas
%% em Sousa2026.

Monteiro et al.~\cite{Monteiro2022} estenderam o ESBMC para C++, com suporte a
\textit{templates} e herança virtual --- diretamente aplicável ao PX4, escrito em
C++17. Wu et al.~\cite{Wu2024} verificaram componentes de arquiteturas Arm
Confidential Computing, demonstrando a escalabilidade da ferramenta em sistemas
complexos. Em outra direção, Menezes et al.~\cite{Menezes2025GCSE}, em
colaboração com a Intel, propuseram a técnica VO-GCSE (\textit{Verification
Optimization through Global Common Subexpression Elimination}), que reduz o tempo
de verificação ao eliminar subexpressões redundantes em modelos simbólicos ---
artigo da trilha de demonstrações do FSE 2025. A otimização não foi empregada
nesta pesquisa, mas é relevante para módulos de firmware com estruturas
repetitivas como as do PX4. Na fronteira entre verificação e teste, Alshmrany et
al.~\cite{Alshmrany2022} propuseram o FuSeBMC, que combina BMC e \textit{fuzzing}
para ampliar a cobertura na geração automática de casos de teste.
%% LC (resolvido): informado que a otimizacao VO-GCSE nao foi usada neste
%% trabalho. LC (resolvido): FuSeBMC retirado deste paragrafo (discutido na
%% subsecao de abordagens nao-BMC, como desenvolvimento relacionado).

\section{Verificação de Sistemas Autônomos}
\label{sec:rel-autonomos}

Wanner et al.~\cite{Wanner2024} revisam a certificação e a confiabilidade de
sistemas aviônicos de UAVs, destacando a necessidade de verificação rigorosa, mas
sem aplicá-la ao código-fonte. Siddiqi et al.~\cite{Siddiqi2022} e Mekdad et
al.~\cite{Mekdad2023} identificam fragilidades em UAVs sob a perspectiva de
ataques externos e de rede, sem abordar o firmware embarcado.
Luckcuck et al.~\cite{Luckcuck2019} levantam, em \textit{survey} sistemático, os
formalismos e desafios da verificação de sistemas robóticos autônomos,
evidenciando que a verificação de firmware de autopiloto permanece pouco
explorada.

Taylor et al.~\cite{Taylor2021} propõem o Avis, um \textit{model checking}
\textit{in-situ} para VANTs que executa o firmware em simulação e injeta falhas
de sensores para observar o comportamento do veículo. É o trabalho mais próximo
quanto ao alvo, porém verifica o sistema em execução, e não o código-fonte por
análise estática, como o BMC aqui adotado. Bernardeschi et
al.~\cite{Bernardeschi2023} combinam verificação formal e co-simulação no
controle cooperativo de drones, com especificações em lógica temporal, mas não
tratam de defeitos no nível dos \textit{drivers}. Pollien et
al.~\cite{Pollien2021} verificam com Frama-C a biblioteca matemática do
autopiloto \textit{open-source} Paparazzi, provando ausência de erros de execução
em funções de conversão de ponto flutuante; a tese que consolida o
trabalho~\cite{Pollien2023} combina Frama-C e o assistente de provas Coq. Trata-se
de verificação dedutiva, que exige anotações manuais extensas, diferentemente do
BMC automatizado adotado aqui.

Di Sorbo et al.~\cite{deSorbo2022} caracterizam automaticamente problemas de
\textit{safety} reportados em plataformas de software para UAVs, revelando
padrões recorrentes de defeitos. Mzid~\cite{Mzid2024} propõe padrões de projeto
para sistemas embarcados de tempo real \textit{safety-critical}, sem aplicação
empírica a controle de voo. Khatiri et al.~\cite{Khatiri2024} propõem o
Aerialist, que automatiza a geração e a execução de cenários de teste no PX4 por
simulação; embora focado em teste, complementa esta pesquisa ao reforçar o valor
de abordagens automatizadas de garantia de qualidade em firmware de autopiloto.

\subsection{Abordagens não Baseadas em \textit{Model Checking}}
\label{subsec:rel-nao-bmc}

As seções anteriores discutem, em sua maior parte, a própria família de técnicas
desta pesquisa. Convém situá-la também em relação às alternativas, porque a
escolha do BMC implica renúncias que só se tornam visíveis por contraste.

\textbf{Verificação dedutiva e análise estática \textit{sound}.} A verificação
dedutiva estabelece a correção por prova, e não por exploração de estados --- é a
abordagem de Pollien et al.~\cite{Pollien2021} sobre o Paparazzi. Quando as
especificações e as hipóteses do modelo são adequadas, essa verificação pode
estabelecer propriedades para além de um limite finito de execução, diferentemente
do BMC convencional, cuja análise está condicionada à profundidade de
desenrolamento; em contrapartida, exige anotar o código com pré-condições,
pós-condições e invariantes; em uma
base de cerca de 710 mil linhas, escrita por centenas de colaboradores e sem
anotações, esse custo é proibitivo. A interpretação abstrata é \textit{sound} por
construção e escala bem, mas produz alarmes que o analista descarta um a um, sem
o contraexemplo concreto que o BMC entrega --- justamente o que permite a
conferência contra o código de produção descrita na
Seção~\ref{subsec:etapa-replay}. O Astrée, de Cousot et
al.~\cite{Cousot2005Astree}, é o representante mais conhecido dessa família, com
histórico em código aeronáutico de produção, como o da Airbus~\cite{Bertrane2010};
a ferramenta comercial Polyspace segue a mesma linha. Já a linguagem SPARK, um
subconjunto de Ada, permite provar estaticamente a ausência de erros de execução
na compilação~\cite{McCormick2015}. As três compartilham o custo estrutural que
motiva a escolha do BMC: exigem disciplina de anotação ou de linguagem que uma
base pré-existente e não anotada, como a do PX4, não satisfaz.

\textbf{Síntese de código verificado.} Em vez de verificar código existente, o
programa HACMS da DARPA construiu o firmware de VANTs a partir de componentes
concebidos para serem verificados --- entre eles o microkernel seL4 e o
compilador CompCert ---, produzindo o autopiloto SMACCMPilot, que resistiu a um
\textit{red team} com seis semanas e acesso total ao código~\cite{Fisher2017}. O
contraste com esta pesquisa é nítido: o HACMS parte de uma base construída para a
verificação, enquanto o objeto deste trabalho é um firmware de produção
pré-existente e não anotado, sobre o qual o BMC é aplicável sem reescrita.

\textbf{Teste dirigido a VANT.} No extremo oposto estão as técnicas dinâmicas.
Kim et al.~\cite{Kim2021PGFuzz} propõem o PGFUZZ, \textit{fuzzing} guiado por
política aplicado a ArduPilot, PX4 e Paparazzi. O RVFuzzer, de T.~Kim et
al.~\cite{Kim2019RVFuzzer}, usa realimentação do controlador para gerar valores
de parâmetros --- enviados pela estação de solo --- que levam o veículo a estados
inseguros, encontrando defeitos de validação de entrada em autopilotos de
produção, incluindo o PX4. Han et al.~\cite{Han2024RangeFuzz} detectam defeitos de
especificação de faixa em módulos de controle de voo. Os três encontram defeitos
reais, por execução, não por análise. A relação com este trabalho é de
complementaridade: técnicas dinâmicas exercitam o sistema integrado, com
escalonador e hardware presentes, mas cobrem apenas as execuções que ocorrem; o
BMC explora exaustivamente os caminhos representados no modelo até o limite de
desenrolamento, sobre um recorte abstraído por \textit{stubs}. O argumento é desenvolvido na
Seção~\ref{sec:discussao-resultados}.

\textbf{Contexto normativo.} Na aviação certificada, a DO-178C rege a garantia de
software aeronáutico crítico~\cite{RTCA_DO178C}, e seu suplemento DO-333
estabelece as condições sob as quais métodos formais podem substituir, total ou
parcialmente, atividades de teste exigidas pela norma~\cite{RTCA_DO333}. É nesse
regime que Astrée, Polyspace e SPARK acumularam histórico de uso. O PX4, por não
ser certificado sob a DO-178C, está fora desse arcabouço: a verificação aqui
conduzida não substitui nem satisfaz requisito de certificação algum, e seus
resultados são evidência técnica independente, não conformidade normativa.

\section{Posicionamento deste Trabalho}
\label{sec:rel-posicionamento}

A Tabela~\ref{tab:trabalhos-relacionados} sintetiza os principais trabalhos
relacionados quanto à técnica, ao sistema-alvo, ao tipo de estudo, ao objeto
avaliado e à ferramenta.

\begin{table}[p]
\centering
\caption{Síntese comparativa dos trabalhos relacionados em verificação formal}
\label{tab:trabalhos-relacionados}
\scriptsize
\renewcommand{\arraystretch}{0.92}
\setlength{\tabcolsep}{2pt}
\begin{tabular}{|p{2.3cm}|p{2.0cm}|p{2.2cm}|p{1.7cm}|p{2.1cm}|p{1.9cm}|}
\hline
\textbf{Trabalho} & \textbf{Técnica} & \textbf{Sistema-alvo} &
\textbf{Tipo de estudo} & \textbf{Objeto} & \textbf{Ferramenta} \\
\hline
Beyer (2022)        & Survey MC       & Múltiplas linguagens    & Survey     & ---                   & Vários    \\ \hline
Beyer (2024)        & SV-COMP         & Benchmark C             & Competição & Benchmark             & Vários    \\ \hline
Kroening (2014)     & BMC             & Programas C             & Ferramenta & Benchmark             & CBMC      \\ \hline
Menezes (2024)      & BMC             & Programas C             & Ferramenta & Benchmark             & ESBMC     \\ \hline
Sousa (2026)        & BMC             & Open-source C           & Empírico   & Sistema de produção   & ESBMC     \\ \hline
Monteiro (2022)     & BMC             & Programas C++           & Ferramenta & Código sintético      & ESBMC     \\ \hline
Wu (2024)           & BMC             & Arm CCA                 & Empírico   & Sistema de produção   & ESBMC     \\ \hline
Alshmrany (2022)    & BMC + Testes    & Programas C             & Ferramenta & Benchmark             & FuSeBMC   \\ \hline
Menezes (2025)      & BMC + GCSE      & Programas C             & Demonstração & Benchmark e indústria & ESBMC   \\ \hline
Fisher (2017)       & Síntese verif.  & \textbf{UAV SMACCMPilot}& Empírico   & Sistema de produção   & seL4/CompCert \\ \hline
Pollien (2021)      & Dedutivo        & UAV Paparazzi           & Empírico   & Sistema de produção   & Frama-C   \\ \hline
Cousot (2005)       & Interp. abstrata& Código aeroespacial     & Ferramenta & Sistema de produção   & Astrée    \\ \hline
McCormick (2015)    & Dedutivo/ling.  & Ada embarcado           & Conceitual & ---                   & SPARK     \\ \hline
Mzid (2024)         & Padrões         & Sistemas RT críticos    & Conceitual & ---                   & --        \\ \hline
Taylor (2021)       & MC in-situ      & \textbf{UAV runtime}    & Empírico   & Sistema de produção   & Avis      \\ \hline
Bernardeschi (2023) & FV + co-sim     & \textbf{Drones cooper.} & Empírico   & Modelo e co-simulação & NuSMV     \\ \hline
Luckcuck (2019)     & Survey FV       & Robótica autônoma       & Survey     & ---                   & Vários    \\ \hline
Wanner (2024)       & Survey          & \textbf{UAV aviônico}   & Survey     & ---                   & --        \\ \hline
Siddiqi (2022)      & Survey          & \textbf{UAV segurança}  & Survey     & ---                   & --        \\ \hline
Mekdad (2023)       & Survey          & \textbf{UAV redes}      & Survey     & ---                   & --        \\ \hline
Kim (2021)          & Fuzzing         & \textbf{PX4, ArduPilot} & Empírico   & Sistema de produção   & PGFUZZ    \\ \hline
Kim (2019)          & Fuzzing         & \textbf{PX4, ArduPilot} & Empírico   & Sistema de produção   & RVFuzzer  \\ \hline
Han (2024)          & Fuzzing         & \textbf{Controle de voo}& Empírico   & Sistema de produção   & ---       \\ \hline
Khatiri (2024)      & Simulação       & \textbf{PX4 VANTs}      & Ferramenta & Simulação             & Aerialist \\ \hline
Di Sorbo (2022)     & Análise SW      & \textbf{UAV software}   & Empírico   & Relatos de defeito    & --        \\ \hline
\textbf{Este trabalho} & \textbf{BMC} & \textbf{PX4 firmware}   & \textbf{Empírico} & \textbf{Código-fonte (1 VAM)$^\dagger$} & \textbf{ESBMC (+ CBMC)} \\ \hline
\end{tabular}

\scriptsize
\begin{flushleft}
$^\dagger$VAM = violação alcançável no modelo, na acepção fixada no
Capítulo~\ref{cap:introducao}. Não se reivindica defeito confirmado: a análise do
driver de produção mostrou o estado violador inalcançável
(Seção~\ref{subsec:gps-inducao}). Ver Tabela~\ref{tab:resultados}.
\end{flushleft}
\end{table}

Este trabalho diferencia-se por aplicar BMC diretamente ao código-fonte de
\textit{drivers} de um sistema de produção real e amplamente utilizado (PX4
Autopilot, ramo \textit{main}, \textit{commit} \texttt{8be2590c15}). Enquanto
Taylor et al.~\cite{Taylor2021} verificam o comportamento em execução e
Bernardeschi et al.~\cite{Bernardeschi2023} focam a lógica de controle
cooperativo, este estudo verifica \textit{drivers} críticos no nível do
código-fonte, identificando candidatos a defeito de memória e aritmética em
caminhos representados no modelo de verificação, cuja exploração não é garantida
pelos testes convencionais do projeto. Em relação ao
fluxo de BMC de Sousa et al.~\cite{Sousa2026}, do qual herda as etapas gerais, a
contribuição metodológica está em dois passos próprios de firmware: a abstração
da E/S de hardware por \textit{stubs} e a especificação de \textit{assertions}
derivadas de \textit{datasheets} e do código-fonte.

Diferentemente dos \textit{surveys} sobre segurança de
UAVs~\cite{Wanner2024, Siddiqi2022, Mekdad2023}, este trabalho oferece evidência
empírica da \textbf{aplicabilidade} da verificação formal a firmware de drones em
produção: uma violação alcançável no modelo (GPS uBlox), reportada à comunidade
PX4 (Issue~\#26866), proposta como endurecimento defensivo (PR~\#27576) e mostrada
inalcançável no driver de produção, e dois candidatos triados, respectivamente,
como especificação disputada (temperatura do BMI088) e falso positivo (FIFO do
BMI088), com o procedimento de triagem documentado. A aplicação ao SAGRES e ao
monitoramento ambiental na Amazônia é a motivação do estudo de caso, e não um
fator que altere o método.
%% LC (resolvido): retirada a sugestao de que o monitoramento ambiental influencia
%% o metodo -- apresentado como motivacao, nao como fator metodologico.

\section{Resumo}
\label{sec:rel-resumo}

Este capítulo posicionou a pesquisa no estado da arte. Ferramentas de BMC como
ESBMC e CBMC têm eficácia crescente na detecção de defeitos em código C/C++ de
produção, mas a maior parte dos trabalhos ainda se concentra em \textit{benchmarks}
sintéticos ou aplicações genéricas. Entre os sistemas autônomos, as abordagens
existentes tratam de verificação em execução, co-simulação ou comportamento
emergente, sem alcançar o nível dos \textit{drivers} de firmware. As alternativas
ao \textit{model checking} --- verificação dedutiva, análise estática
\textit{sound}, síntese de código verificado e teste dirigido --- foram discutidas
na Seção~\ref{subsec:rel-nao-bmc}, e sua relação com esta pesquisa é de
complementaridade.

O trabalho de Pollien et al.~\cite{Pollien2021} sobre o Paparazzi é o precedente
mais próximo em verificação formal de código de autopiloto, mas usa técnicas
dedutivas com anotações manuais, ao passo que esta pesquisa emprega BMC
automatizado com geração de contraexemplos concretos. A combinação de verificação
automatizada, conferência contra o código de produção, validação cruzada com o
CBMC e reporte direto à comunidade \textit{open-source} é a abordagem proposta
para firmware de autopiloto de VANTs, endereçando uma lacuna observada no
levantamento.
